% Generated draft: requires mathematical and source review.
\documentclass[11pt]{article}
\usepackage[margin=25mm]{geometry}
\usepackage{fontspec}
\setmainfont{Latin Modern Roman}
\usepackage{amsmath,amssymb,mathtools,unicode-math}
\setmathfont{Latin Modern Math}
\usepackage{enumitem,needspace,xurl,hyperref}
\setlength{\parindent}{0pt}
\setlength{\parskip}{5pt}
\setlength{\emergencystretch}{3em}
\newcommand{\N}{\mathbb{N}}
\newcommand{\Z}{\mathbb{Z}}
\newcommand{\Q}{\mathbb{Q}}
\newcommand{\R}{\mathbb{R}}
\newcommand{\C}{\mathbb{C}}
\newcommand{\F}{\mathbb{F}}
\newcommand{\A}{\mathbb{A}}
\newcommand{\PP}{\mathbb{P}}
\newcommand{\E}{\mathbb{E}}
\newcommand{\eps}{\varepsilon}
\newcommand{\dd}{\,\mathrm{d}}
\begin{document}
\section*{UnsolvedMath ID: 30001487}
\textbf{Source problem code:} catalogue entry
\subsection{Definitions and mathematical statement}
An infinite-dimensional Lie group is a group equipped with an infinite-dimensional smooth manifold structure for which multiplication and inversion are smooth, relative to a specified category of manifolds and a specified notion of smoothness. The supplied record does not specify these choices.

In a common groupoid formulation, a Lie \(2\)-group is a smooth groupoid equipped with a monoidal structure, coherent associativity and unit isomorphisms, and inverses for objects up to isomorphism. Strict formulations impose corresponding identities strictly. These descriptions explain the terminology; the source does not select a strict or weak formulation or an ambient smooth category.

The supplied research question is:
\begin{quote}
What can be learned from infinite-dimensional Lie groups about Lie \(2\)-groups, and conversely?
\end{quote}
Both directions of the question are retained. The record specifies neither classes of objects to quantify over nor a particular construction, comparison, or conclusion to establish. Consequently, it determines a research direction but does not determine a precise quantified conjecture. In particular, no equivalence or correspondence theorem is asserted by this draft.
\subsection{Short English statement}
The question asks how the theories of infinite-dimensional Lie groups and Lie \(2\)-groups can inform one another. The intended smooth categories, strictness conventions, and mathematical targets remain unspecified, so a precise theorem cannot be extracted without additional information.
\subsection{Sources}
\begin{itemize}
\item Supplied local UnsolvedMath registry record: permanent ID \(30001487\), exact source problem code \texttt{catalogue entry}, in the Oberwolfach Reports Open Problems collection. The record assigns the year \(2010\); no individual proposer is named.
\item Supplied report citation: \emph{Geometry, Quantum Fields, and Strings: Categorial Aspects} (2010), Oberwolfach Reports, pp.\ 1485--1533, DOI \texttt{10.4171/owr/2010/25}. Supplied source URL and literature reference labeled OWR 2010/25: \url{https://doi.org/10.4171/owr/2010/25}. The supplied background describes these reports as documenting mathematical workshops and their open-problem discussions.
\item Supplied literature assessment, checked \texttt{2026-08-21}: the record labels the question open and states that integration and higher-gauge-group results address special cases without providing a general resolution. This is supplied dated curation, not an independently verified assessment. The background explicitly calls for independent verification of time-sensitive status and cited claims.
\item The registry labels the statement exact and reports that manual review found the extracted wording self-contained, mathematically coherent, and free of LaTeX or report-layout artifacts. These are supplied verification claims; the report has not been independently consulted for this draft.
\end{itemize}
\end{document}
